\chapter{Воспроизводимость и трассируемость}
\label{app:reproducibility}

\section{Назначение приложения}

Воспроизводимость в этой работе разделяется на два уровня. Первый уровень ---
\emph{воспроизводимость документа}: из конкретного состояния Git-репозитория
проверяются контракты данных, предназначенные для диссертации,
детерминированно формируются таблицы LaTeX и собирается PDF. Второй уровень ---
\emph{научная трассируемость}: утверждения диссертации связываются с
отслеживаемыми научными сводками либо с замороженными идентификаторами
доказательных материалов, целостность которых была зафиксирована. Эти два
уровня не следует смешивать: повторная сборка PDF не является повторным
выполнением научного эксперимента.

\section{Канонический путь сборки документа}

Основная точка входа для сборки в репозитории:

\begin{verbatim}
make thesis
\end{verbatim}

Перед XeLaTeX выполняется \code{scripts/build_thesis_assets.py}. Скрипт
проверяет контракт утверждений, сводку QWake C2 для диссертации, привязки
SHA-256 к отслеживаемым доказательным материалам этап~1/2/3 и выбранные
численные агрегаты. Только после успешной проверки создаются файлы в
\code{thesis/generated/}. Эти сгенерированные файлы не фиксируются в Git и не
являются самостоятельными научными доказательными материалами.

Проверка контрактов данных без установленного LaTeX выполняется командой:

\begin{verbatim}
make thesis-check
\end{verbatim}

Каноническая сборка CI дополнительно фиксирует коммит и дерево Git, версии
\texttt{latexmk}, XeLaTeX и Biber, проверяет сигнатуру и размер PDF и публикует
SHA-256 конкретного артефакта PDF.

\section{Слои доказательных материалов для диссертации}

В документе используются два компактных машиночитаемых слоя:
\begin{itemize}
  \item \code{thesis/data/core_results_verified_summary.json} --- проекция
    результатов этап~1/2/3. Для каждого исходного отслеживаемого артефакта
    хранится ожидаемый SHA-256, который повторно проверяется при сборке;
  \item \code{thesis/data/qwake_c2_verified_summary.json} --- компактное
    представление результата QWake C2 с зафиксированной целостностью и
    независимо проверенных агрегатов. Сборка дополнительно сверяет статусы
    C08--C11 с замороженным результатом отбора, агрегатами лучшего безопасного
    правила и границей протокола, не выполняя новую оценку правил и не изменяя
    модель учёта стоимости. Файл явно маркирован значением
    \code{not_new_scientific_evidence}.
\end{itemize}

Полный набор проверочных материалов QWake намеренно не копируется в обычное
дерево исходных файлов диссертации. Вместо этого сохраняются замороженные
идентификаторы артефактов, квитанций, выбранного правила и общей
последовательности учёта стоимости принятия решения. Это позволяет проверить,
к какому именно результату с зафиксированной целостностью относится текст, не
создавая ложного впечатления, что компактный JSON для диссертации заменяет
исходный набор научных материалов.

\section{Связанные источники и контрольные суммы}

Следующая таблица формируется непосредственно из контрактов данных для
диссертации во время сборки. Для этап~1/2/3 она перечисляет исходные
отслеживаемые артефакты и их проверяемые SHA-256. Для QWake она перечисляет
замороженные идентификаторы, которыми ограничен итоговый результат C2.

\input{generated/reproducibility_manifest}

\section{Границы протокола QWake}

Для QWake C2 зафиксированы следующие факты протокола: правило C2 не было
зафиксировано, C3 не открыт, повторная попытка и повторное использование
разрешения не допускаются. Поэтому воспроизводимость диссертации не включает
повторное выполнение C2 или переход к C3. Любой будущий эксперимент по
добавочной стоимости должен иметь новую оцениваемую величину, новый протокол и
новую цепочку доказательных материалов; он не является продолжением уже
закрытого вызова C2.

\section{Минимальная проверка полученного PDF}

После успешной сборки конкретный артефакт может быть идентифицирован командами:

\begin{verbatim}
pdfinfo thesis/main.pdf
sha256sum thesis/main.pdf
\end{verbatim}

SHA-256 в этом контексте идентифицирует конкретный собранный PDF. Побайтовая
идентичность PDF между независимыми сборками текущим процессом сборки не
заявляется как отдельный научный результат. Источником содержательной
воспроизводимости остаётся точное состояние исходного кода Git вместе с
успешными проверками контрактов данных и происхождения артефактов.

\section{Практическая независимая проверка}

Для независимого чтения и проверки документа достаточно следующей
последовательности:
\begin{enumerate}
  \item получить конкретный коммит Git;
  \item выполнить \texttt{make thesis-check} и убедиться, что проверки
    утверждений, арифметики и границ протокола QWake, а также привязок основных
    результатов проходят успешно;
  \item выполнить используемую в CI проверку эпистемического языка;
  \item в среде с XeLaTeX/Biber выполнить \texttt{make thesis};
  \item проверить метаданные PDF и SHA-256;
  \item для спорного численного утверждения перейти от сгенерированной таблицы
    к JSON для диссертации, а затем к указанному исходному артефакту или
    замороженному научному идентификатору.
\end{enumerate}

Такой порядок делает путь от текста к доказательным материалам явным и
одновременно сохраняет границу между документированием уже завершённого
исследования и новым научным выполнением.
